\subsection{IRREDUCIBILITY OF THE GEOMETRIC REPRESENTATION OF A COXETER GROUP}

We retain the notation of the preceding nos., and assume that S is finite.

PROPOSITION 7. Assume that (W,S) is irreducible (Chap. IV, §1, no. 9). Let \(\mathbf{E}^0\) be the subspace of E orthogonal to E with respect to \(\mathbf{B}_{\mathbf{M}}\) . The group W acts trivially on \(\mathbf{E}^0\) , and every subspace of E distinct from E and stable under W is contained in \(\mathbf{E}^0\) .

If \(x \in \mathrm{E}^0\) , then \(\sigma_s(x) = x - 2\mathrm{B}_{\mathrm{M}}(e_s, x)e_s = x\) for all \(s \in \mathrm{S}\) . Since \(\mathrm{W}\) is generated by \(\mathrm{S}\) , it follows that \(\mathrm{W}\) acts trivially on \(\mathrm{E}^0\) .

Let \(E'\) be a subspace of \(E\) stable under \(W\) . Let \(s, s' \in S\) be two elements that are linked in the graph \(\Gamma\) of \((W, S)\) (Chap. IV, §1, no. 9); recall that this means that \(m(s, s') \geqslant 3\) . Suppose that \(e_s \in E'\) . Then \(\sigma_{s'}(e_s) \in E'\) and since the coefficient of \(e_{s'}\) in \(\sigma_{s'}(e_s)\) is non-zero, we have \(e_{s'} \in E'\) . Since \(\Gamma\) is connected, it follows that, if \(E'\) contains one of the \(e_s\) it contains them all and coincides with \(E\) . Except in this case, it follows from Prop. 3 of §2, no. 2 that, for all \(s \in S\) , \(E'\) is contained in the hyperplane \(H_s\) orthogonal to \(e_s\) . Since the intersection of the \(H_s\) is equal to \(E^0\) , this proves the proposition.

COROLLARY. Assume that (W, S) is irreducible. Then:

\begin{enumerate}
\def\labelenumi{\alph{enumi})}
\item
  If \(\mathrm{B_M}\) is non-degenerate, the W-module E is absolutely simple.
\item
  If \(\mathrm{B_M}\) is degenerate, the W-module E is not semi-simple.
\end{enumerate}

In case \(a\) ), Prop. 7 shows that \(E\) is simple, hence also absolutely simple (§ 2, no. 1, Prop. 1).

In case \(b\) ), \(\mathrm{E}^0 \neq 0\) , \(\mathrm{E} \neq \mathrm{E}^0\) (since \(\mathrm{B}_{\mathrm{M}} \neq 0\) ), and Prop. 7 shows that \(\mathrm{E}^0\) has no complement stable under W; thus, the W-module E is not semi-simple.

